% =====================================================================
%  一、问题重述
% =====================================================================
\section{问题重述}

\subsection{问题背景}
无线电测向与干扰源定位是频谱管理、电磁环境监测中的重要任务。借助可移动测向机，
通过在不同检测点测得干扰源的示向度（方位角），即可对干扰源位置进行交会估计；
在此基础上进一步规划检测路线、逐一清除干扰源。\textbf{【待填充：结合赛题原文，
补充背景与任务场景，注意不要照抄原题，用自己语言概述。】}

\subsection{问题的提出}
根据赛题，本文需要依次解决以下四个子问题：

\begin{enumerate}
    \item \textbf{问题一}：……在已知两个检测点及其示向度（含测量误差）的情况下，
          求交会定位区域的直径，并判断以该直径为直径的圆能否覆盖整个定位区域。
          \textbf{【待填充：按原题准确复述。】}

    \item \textbf{问题二}：……在给定第一检测点测得示向度后，研究第二检测点的
          最优选择策略，并给出相应的候选区域。
          \textbf{【待填充：按原题准确复述。】}

    \item \textbf{问题三}：……在半径为 $1800\,\mathrm{m}$ 的圆形区域内存有
          $10$--$16$ 个位置未知的全向干扰源，要求在规定现实时间（$20$ 分钟）内
          将全部干扰源定位并清除。
          \textbf{【待填充：按原题准确复述，含测向、换频道、移动、清除等时间约束。】}

    \item \textbf{问题四}：……在问题三基础上叠加若干方向未知的定向干扰源，
          要求仍能在限定时间内完成定位与清除。
          \textbf{【待填充：按原题准确复述。】}
\end{enumerate}
